\section{The original projective tuple map on a unit coordinate}

\begin{definition}[The same original graded tuple evaluation]
\label{def:fc-generated-tuple-graded-map}
\leanok
\lean{EnvelopeProperty.GeneratedTupleProjectiveCoordinates.tupleGradedMap, EnvelopeProperty.GeneratedTupleProjectiveCoordinates.tupleGradedMap_X, EnvelopeProperty.GeneratedTupleProjectiveCoordinates.tupleGradedMap_irrelevant, EnvelopeProperty.GeneratedTupleProjectiveCoordinates.affineTupleMap_graded}
For an actual unit-ideal tuple $(a_i)$ in the original domain $A$ and
an actual coefficient map $R\to A$, compose the original polynomial
coefficient map with the original degree-one Rees evaluation. This
graded map sends every original $X_i$ to the same $a_iT$.
\end{definition}
\begin{proof}
Both maps preserve the original grading. Their composite has the
required irrelevant-ideal containment by composition of the genuine
coefficient and surjective Rees-map containments. Evaluation on an
original variable is the defining degree-one Rees evaluation.
\end{proof}

\begin{theorem}[Exact original coordinate chart of the same tuple morphism]
\label{thm:fc-generated-tuple-unit-chart}
\leanok
\lean{EnvelopeProperty.GeneratedTupleProjectiveCoordinates.chartRingMap, EnvelopeProperty.GeneratedTupleProjectiveCoordinates.affineTupleMap_coordinateChart, EnvelopeProperty.GeneratedTupleProjectiveCoordinates.chartRingMap_coordinate, EnvelopeProperty.GeneratedTupleProjectiveCoordinates.chartRingMap_C, EnvelopeProperty.GeneratedTupleProjectiveCoordinates.chartRingMap_eq_eval, EnvelopeProperty.GeneratedTupleProjectiveCoordinates.affineTupleMap_eq_eval_chart}
\uses{def:fc-generated-tuple-graded-map,thm:fc-original-unit-rees-chart-section,thm:fc-original-unit-rees-coordinate-ratio}
If the original $a_i=1$, the actual original tuple morphism is
the actual standard projective chart inclusion composed with the
spectrum of polynomial evaluation $X_j/X_i\mapsto a_j$, with the
original coefficient map. In particular this is a whole morphism
identity for the same original Rees-constructed tuple morphism.
\end{theorem}
\begin{proof}
The original homogeneous Rees denominator, namely the image of
$X_i$, is precisely $T$. The genuine unit-blowup inverse is the
actual homogeneous chart section from the preceding result. Compose
with the native Proj chart square for the original composite graded
map. On each original coordinate fraction its coefficient is
$(a_jT)/T=a_j$; on original constants it is the original coefficient
map. Polynomial-ring extensionality identifies the entire chart
homomorphism with this evaluation. The original homogeneous-chart
isomorphism then gives the stated whole projective morphism identity.
\end{proof}
